% Isolated after the CP27 canonical freeze; not part of that frozen paper.
\subsection{The terminal Frobenius pair in the quotient induction}

Retain the hypotheses and notation of B74: \(p\ge5\),
\(q=p^2\), \(N=q-1\), \(\mu_p\subset K\), and the
absolute ramification \(e\) is prime to \(p\).
Local arguments use algebraically closed residue field and
the actual common good tame field of degree \(N\).

\begin{lemma}[The enlarged terminal interval]
\label{lem:B-RD75-terminal}
\status{Candidate; complete proof supplied, independent audits pending.}
Suppose two original two-slope groups have parameters
\(B_i=pb_i\), common high-line character, and common finite
determinant. If \(1\le b_i\le p\), their finite special
inertia characters agree.
\end{lemma}
\begin{proof}
The same character equations as (B-RD74.4) give
\[
 b_1\equiv b_2\pmod{p-1},\qquad
 c_2-c_1\equiv p(b_2-b_1)\pmod N,
 \qquad c_i\equiv pb_i\pmod{p+1}.
 \tag{B-RD75.1}
\]
Equal \(b_i\) give equal \(c_i\pmod N\). The only unequal
congruent pair in \(\{1,\ldots,p\}\) is \(\{1,p\}\).
Reorder to \(b_1=1,b_2=p\). Then
\[
 c_1\equiv p\pmod{p+1},\qquad
 c_2-c_1\equiv p(p-1)\equiv(p-1)c_1\pmod N.
 \tag{B-RD75.2}
\]
Thus \(c_2\equiv pc_1\pmod N\), and
\(\{c_2,pc_2\}=\{c_1,pc_1\}\) because \(p^2\equiv1\pmod N\).
These are precisely the finite special eigenvalue multisets.
The parameters need not be equal: only the requested character
is being recovered. This calculation is at the original
marked tangent exponents, so no change of the common inertia
generator or covariant convention is implicit.
\end{proof}

\begin{proposition}[A larger induction upper range]
\label{prop:B-RD75-formal}
\status{Candidate; complete proof supplied, independent audits pending.}
Let \(n\ge3\). Whole potentially supersingular height-two
groups have their finite special characters determined by
actual \(p^n\)-torsion whenever
\[
 e<p^{n-1}(p+1).
 \tag{B-RD75.3}
\]
\end{proposition}
\begin{proof}
Run the complete actual-quotient induction of
Proposition~\ref{prop:B-RD74-formal} unchanged through all
its stopping cases. The actual quotient level, semilinear
descent, high-line recovery, prime-parameter case and tame/wild
case do not use its upper bound on \(e\). That bound is used
only at the final all-small stage. Under the present bound,
that stage instead gives
\[
 0<b_i<\frac{e}{p^{n-1}}<p+1,
 \qquad b_i\in\{1,\ldots,p\}.
 \tag{B-RD75.4}
\]
The preceding lemma recovers the original finite character
also in the additional pair \(\{1,p\}\). Every other branch
was already settled before this terminal interval was used.
The full proof still retains only actual \(p\)-torsion after
the last quotient; no further quotient is taken.
\end{proof}

\begin{theorem}[A longer third-index interval]
\label{thm:B-RD75}
\status{Candidate; complete proof supplied, independent audits pending.}
In the full B.4 height-two monodromy setting with perfect
residue field, \(p\ge5\), \(\mu_p\subset K\), and \(p\nmid e\),
every \(n\ge4\) satisfies
\[
 e<p^{n-1}(p+1)\quad\Longrightarrow\quad
 \nu_{\min}(p,e,2)\le n-1.
 \tag{B-RD75.5}
\]
Consequently
\[
 p^3<e<p^3(p+1)\quad\Longrightarrow\quad
 \boxed{\nu_{\min}(p,e,2)=3}.
 \tag{B-RD75.6}
\]
For \(p=5\), every permitted multiple of four from \(128\)
through \(748\), prime to five, therefore has exact index three.
\end{theorem}
\begin{proof}
The determinant and height-one filtered reductions in
Theorem~\ref{thm:B-RD74-upper} are unchanged. Its mixed
whole/filtered comparison is excluded because, for \(n\ge4\),
\[
 e<p^{n-1}(p+1)<p^{2n-3}.
 \tag{B-RD75.7}
\]
The second strict inequality is equivalent to
\(p+1<p^{n-2}\), valid for \(p\ge5,n\ge4\).
The preceding proposition supplies the whole-group upper
bound, completing the full-monodromy argument.
Taking \(n=4\) and using the B63 actual \(p^3\)-torsion
lower construction whenever \(e>p^3\) makes the two bounds
equal to three. At \(p=5\), the interval is \((125,750)\),
whose first and last permitted multiples of four are \(128\)
and \(748\).
\end{proof}

\begin{remark}[The first nontrivial terminal collision remains open]
\label{rem:B-RD75-scope}
\status{Candidate scope; no exact claim beyond the stated interval.}
The extra terminal pair is harmless because it preserves the
unordered finite character; no formal-group existence is
inferred from its arithmetic compatibility. At the next
terminal value \(p+1\), the pair \(\{2,p+1\}\) need not be
a Frobenius orbit. This is an unresolved arithmetic possibility,
not a constructed comparison of actual torsion groups.
For \(p=5,e=752\), the proved bounds would still be
\(3\le\nu_{\min}\le4\). The general exact third-index range,
wild absolute ramification and higher heights remain open.
A finite exact character/threshold supplement is being prepared;
the argument itself is the all-prime character calculation above.
\end{remark}
